\chapter{Frontend y plug-and-play}
\label{cap:plug-and-play}

\begin{intuicion}
La regla \emph{plug-and-play} dice: cada optimización S/M/L se
activa por defecto al levantar el binario. Sin \emph{flags} a setear,
sin \emph{config} a tocar, sin \emph{toggles} en la UI. El usuario
levanta \texttt{npm run tauri dev} y todo lo que está \emph{shipped}
se nota inmediatamente. Este capítulo documenta cómo se materializa
esa promesa.
\end{intuicion}

\section{Flags default-on al \emph{startup}}

\filepath{apps/tauri/src-tauri/src/lib.rs} y
\filepath{apps/cli/src/main.rs} corren un loop al inicio que setea
flags si el usuario no las exportó:

\begin{lstlisting}[language=Rust]
for (key, default) in [
    ("GRAPHHUNTER_LFTJ_PLANNER", "1"),  // M1 structural K_4
    ("GRAPHHUNTER_K4_LFTJ", "1"),        // S3 name marker
    ("GRAPHHUNTER_YANNAKAKIS", "1"),     // L1 stage 2/3
] {
    if std::env::var(key).is_err() {
        unsafe { std::env::set_var(key, default) };
    }
}
\end{lstlisting}

Override con \texttt{KEY=0} sigue funcionando para usuarios que
quieran el \emph{path} legacy (debugging, A/B comparison).

\section{La pildora del toolbar}

\begin{figure}[H]
\centering
\begin{tikzpicture}[
    node distance=0.4cm,
    every node/.style={font=\scriptsize},
    pill/.style={draw, rounded corners=4mm, minimum width=2.0cm,
                 minimum height=0.7cm, font=\footnotesize}
]
\node[pill, draw=accent, fill=accentlight] (perf) {Perf: M+L1};
\node[pill, draw=green!50!black, fill=green!10, right=of perf]
    (memok) {Mem: 14\%};
\node[pill, draw=orange!70!black, fill=orange!10, right=of memok]
    (memwarn) {Mem: 78\%};
\node[pill, draw=red!70!black, fill=red!10, right=of memwarn]
    (memcrit) {Mem: 92\%};
\end{tikzpicture}
\caption{Pildoras del \emph{toolbar}. La de \emph{perf} indica si
los \emph{flags} M+L1 están activos (verde) o si el usuario forzó el
\emph{legacy path} (gris). La de \emph{mem} cambia de verde a ámbar
(\(\geq 75\%\) RSS/total) a rojo (\(\geq 90\%\)).}
\label{fig:plug-pills}
\end{figure}

\subsection{Implementación}

\filepath{apps/tauri/src-tauri/src/commands/perf.rs} expone
\texttt{cmd\_get\_perf\_status} que retorna una struct con todos los
flags resueltos. \filepath{apps/tauri/src-tauri/src/commands/memory.rs}
expone \texttt{cmd\_get\_memory\_status} con RSS + host total +
ratio.

\filepath{apps/tauri/src/App.tsx} llama ambos comandos al \emph{mount}
y polea el de memoria cada 5~s. Las pildoras del \emph{toolbar} se
renderizan con \emph{data-attributes} que \texttt{App.css} mapea a
los tres estados.

\section{Telemetría adicional en \emph{batch}}

El endpoint \texttt{run\_hunt\_batch} (L2) retorna una respuesta
con tres \emph{cache counters} que el \emph{frontend} puede mostrar
como \emph{tooltip} sobre la pildora de \emph{perf}:

\begin{lstlisting}[language=Rust]
pub struct RunHuntBatchResponse {
    pub entries: Vec<RunHuntBatchEntry>,
    pub indexes_were_warm: bool,
    pub bag_cache_hits: usize,         // L2 stage 2
    pub bag_cache_misses: usize,
    pub result_cache_hits: usize,      // L2 stage 3a
    pub last_prefix_resume_depth: usize,  // L2 stage 3b
    pub prefix_cache_hits: usize,
}
\end{lstlisting}

\section{Doc: ¿qué tocó el usuario?}

La promesa \emph{plug-and-play} se hace verificable: si después de
una optimización el usuario tuvo que cambiar algo (env var, config
file, \emph{toggle}), la promesa se rompió. La regla quedó codificada
en el \emph{memory feedback} \filepath{feedback\_perf\_plug\_and\_play.md}:

\begin{enumerate}
    \item \textbf{Default-on al \emph{startup}}: cada \emph{flag}
    nuevo se setea en \texttt{lib.rs} y \texttt{main.rs} con
    \texttt{if std::env::var(key).is\_err()}.
    \item \textbf{Surface en UI}: extender
    \texttt{cmd\_get\_perf\_status} y la pildora del toolbar.
    \item \textbf{Default-off sólo con causa}: \emph{flags} que
    reescriben el DFS \emph{hot loop} (AMAC, \emph{semi-join})
    quedan opt-in durante un \emph{soak window}; cuando los
    \emph{parity tests} aguantan \(\geq 1\) semana, se flipean.
    \item \textbf{Doc obliga}: cada \emph{flip} actualiza
    \filepath{08-env-vars.tex}, \filepath{12-preguntas-frecuentes.tex}
    y \texttt{make all} para refrescar los 4 PDFs.
    \item \textbf{Capabilities de \emph{build}}: \emph{Cargo features}
    como \texttt{symmetry-nauty} se reflejan via constantes
    \texttt{HAS\_NAUTY}-style en \texttt{cmd\_get\_perf\_status}.
\end{enumerate}

\section{El cierre del libro}

Esta es la última pieza pedagógica del diseño. El resto de los libros
GraphHunter (\emph{Arquitectura}, \emph{Especificación}, \emph{Manual},
\emph{Paper}) profundizan en aspectos específicos:

\begin{itemize}
    \item \emph{Arquitectura}: el \emph{stack} de \emph{crates}, las
    capas, las decisiones de \emph{layering}.
    \item \emph{Especificación}: los \emph{contracts} de HTTP, MCP,
    FFI, las variables de entorno, los formatos de \emph{log} aceptados.
    \item \emph{Manual}: el \emph{user manual} para el analista ---
    cómo se ingiere un dataset, cómo se escribe una hipótesis, cómo se
    interpretan los resultados.
    \item \emph{Paper}: la justificación académica con
    \emph{benchmarks}, \emph{property tests}, y comparaciones contra
    el estado del arte.
\end{itemize}

Este libro es la guía de diseño que los une.
